\section{Conclusion}
This work shows that bio-based particulate-reinforced thermosets can be
 successfully optimised with respect to their compressive yield strength and
 Young's modulus by adapting their composition during manufacture.
Although filler content proved to be the means by which a large degree of
 strength was obtained, careful trade-off between unsaturated polyester 
 and reactive diluent content was necessary in the high-filler region of the
 parameter space to obtain the globally optimal results.
The detailed experiments show how the multi-objective Bayesian optimisation
 employing the log noisy expected hypervolume improvement acquisition
 function was capable of navigating this trade-off, producing a final Pareto
 front consisting of a single dominating point exhibiting
 \(E=5,458\)~\unit{\mega\pascal} and \(\sigma_{y}=165\)~\unit{\mega\pascal}.

In the domain of a sample-constrained single-objective Bayesian optimisation
 using the log noisy expected improvement acquisition function,
 a bio-based composite was obtained which exhibited 69\% of the Young's
 modulus of a conventional styrene-based material which had been optimised
 in the same fashion.
Cross-comparison of the aforementioned multi-objective putative global
 optimum sample with this styrene-based sample developed, showed that it
 attained 75\% the conventional material's strength.

Given the interesting objective topographies witnessed within the parameter
 space with respect to predictor \(x_{3}\) which concerns the balance between
 unsaturated polyester and reactive diluent, future work could examine how
 reactive diluents impact the mechanical properties of these
 materials.
An optimisation of the styrene-based analogues in the single-objective setting
 with respect to the objective of yield strength is lacking herein and would
 be welcomed.